\chapter{State of the Art}
\label{ch:related_work}

% ─────────────────────────────────────────────────────────────────────────────

\section{Generative Models}
\label{sec:generative_models}

\subsection{Generative Adversarial Networks}
\label{subsec:gans}

Generative Adversarial Networks~\cite{goodfellow2014}, proposed by Goodfellow et al.\ in
2014, establish a two-player adversarial game between a generator network $G$ and a
discriminator network $D$. The generator maps samples from a fixed prior distribution
$z \sim p_z$ to the image space; the discriminator receives both real images drawn from
the data distribution $p_{\text{data}}$ and synthetic images produced by $G$, and
attempts to assign them to the correct class. Both networks are trained simultaneously
through the minimax objective:

\begin{equation}
  \min_G \max_D \;\mathbb{E}_{x \sim p_{\text{data}}}[\log D(x)]
  + \mathbb{E}_{z \sim p_z}[\log(1 - D(G(z)))]
  \label{eq:gan_objective}
\end{equation}

At Nash equilibrium, the generator produces images whose distribution is
indistinguishable from the real data distribution. In practice, training is unstable:
the discriminator can saturate early, and the generator can collapse to a small number
of modes. Significant engineering effort, including batch normalisation, spectral
normalisation, and progressive growing, was required before GAN outputs reached
photorealistic quality.

StyleGAN~\cite{karras2019}, introduced by Karras et al.\ in 2019, addresses the quality
and controllability problem through a style-based generator architecture. A mapping
network $f: \mathcal{Z} \rightarrow \mathcal{W}$ transforms the input latent code $z$
into an intermediate style vector $w$, which controls the synthesis network at each
resolution level via adaptive instance normalisation. The spatial content of the image
is separated from its style, and synthesis proceeds by progressively adding detail at
increasing resolutions from $4\times4$ to $1024\times1024$. StyleGAN and its successors
produce human portrait images at a resolution and perceptual quality where gallery-scale
synthetic face collections are visually indistinguishable from stock photography. Dataset~C,
one of the three datasets used in this thesis, is built from StyleGAN-generated faces
paired with FFHQ real portraits.

The instability of GAN training had a secondary effect that is relevant to detection. Early
GAN generators, particularly those based on transposed convolution with fixed integer
strides, converged to outputs with strong periodic artefacts. A transposed convolution with
stride~2 introduces a checkerboard pattern in the spatial domain that appears as spectral
peaks at Nyquist/2 in the FFT magnitude spectrum. These peaks are absent from natural
photographs, whose spatial statistics follow an approximate $1/f$ decay. Detectors that
operated on the frequency spectrum of early GAN outputs could reliably identify these peaks
as evidence of generation. The artefacts were, in effect, a consequence of the discriminator
failing to fully eliminate the upsampling grid structure from the generator's outputs.

Early CNN-based detection methods were designed with GAN outputs in mind. GAN generators
apply learned upsampling operations to increase spatial resolution at each stage; these
operations introduce periodic artefacts in the frequency domain that are absent from
camera-captured images. Detection methods that exploited these spectral signatures worked
well when trained and tested on the same generator family, but the spectral patterns differ
substantially between architectures. A detector trained on ProGAN artefacts does not
generalise to StyleGAN artefacts, because the two architectures use different upsampling
operations with different spectral footprints.

\subsection{Diffusion Models}
\label{subsec:diffusion_models}

Denoising Diffusion Probabilistic Models~\cite{ho2020}, introduced by Ho et al.\ in 2020,
take a different approach to generative modelling. The forward process defines a Markov
chain that gradually corrupts a clean image $x_0$ into isotropic Gaussian noise over $T$
steps:

\begin{equation}
  x_t = \sqrt{\bar{\alpha}_t}\, x_0 + \sqrt{1 - \bar{\alpha}_t}\, \epsilon,
  \qquad \epsilon \sim \mathcal{N}(0, I)
  \label{eq:ddpm_forward}
\end{equation}

where $\bar{\alpha}_t = \prod_{s=1}^t \alpha_s$ is the cumulative noise schedule. A
neural network $\epsilon_\theta(x_t, t)$, typically a U-Net, is trained to predict the
noise $\epsilon$ added at each step. At inference time, generation proceeds by sampling
$x_T \sim \mathcal{N}(0, I)$ and iteratively applying the learned denoising step for
$T$ steps. Because the model is trained on the denoising objective rather than an
adversarial game, training is stable and does not suffer from mode collapse.

Latent Diffusion Models~\cite{rombach2022}, introduced by Rombach et al.\ in 2022, make
this approach computationally tractable at high resolution. An autoencoder compresses
images into a lower-dimensional latent representation; the diffusion process runs in this
latent space rather than in pixel space. A cross-attention mechanism in the denoising
network conditions generation on an external signal, typically a CLIP text embedding,
enabling text-to-image generation. The resulting Stable Diffusion models, beginning with
version~1.5 and progressing through SD~2.1, SDXL, and SD~3, produce high-resolution
photorealistic images from arbitrary text prompts at a computational cost accessible to
consumer hardware. Commercial systems including DALL-E~3 and MidJourney~v6 build on the
same latent diffusion principles. Dataset~A in this thesis uses Stable Diffusion~v1.5
outputs; Dataset~B uses outputs from five diffusion-based generators including SD~2.1,
SDXL, SD~3, DALL-E~3, and MidJourney~v6.

The rise of diffusion models has made detection substantially harder than in the GAN era.
A GAN generator is a fixed mapping from latent space to image space, and the artefacts it
produces are characteristic of that mapping. A diffusion model iteratively refines an
image across hundreds of denoising steps, with each step contributing small corrections
that progressively remove spectral inconsistencies. Modern diffusion outputs at full
resolution are consistently judged as photorealistic by human observers, and the
model-specific artefacts are smaller in magnitude and more variable in character than
those of earlier GAN architectures. Detection methods that relied on generator-specific
spectral signatures face a harder target when the target is a diffusion model rather than
a GAN.

% ─────────────────────────────────────────────────────────────────────────────

\section{AI-Generated Image Detection}
\label{sec:detection_methods}

\subsection{Early Approaches}
\label{subsec:early_detection}

The history of AI image detection mirrors the broader pattern of adversarial dynamics in
computer vision: each generation of detectors is designed around the artefacts of the
generators available at the time, and each generation of generators removes or weakens
those artefacts. Early GAN architectures (2014--2018) produced outputs with strong
checkerboard artefacts from transposed convolution upsampling, visible blurring at
high spatial frequencies, and inconsistent tooth and hair textures. Detectors from the
same period exploited these patterns effectively. The introduction of progressive growing
(2018) and style-based synthesis (2019) produced outputs with substantially weaker
spectral artefacts and higher perceptual quality, breaking earlier detectors. The shift
to diffusion models (2020--2022) removed the fixed upsampling grid structure entirely,
replacing it with iterative denoising that distributes synthesis artefacts over hundreds
of steps; this change rendered spectral detectors unreliable as a class.

Before CNN-based classifiers, AI image detection relied on hand-crafted features drawn
from the signal processing and steganalysis literature. Spatial Rich Model features,
co-occurrence matrices of noise residuals, and JPEG quantisation statistics were used to
characterise the statistical fingerprint of synthesis operations. These methods achieved
reasonable performance on early GAN architectures whose upsampling operations introduced
strong and consistent spectral artefacts. As generator quality improved and upsampling
pipelines became more sophisticated, the artefacts these methods targeted became weaker
and less consistent across architectures, and performance degraded. The shift to CNN-based
feature extraction was motivated by the need for representations that could adapt to the
diversity of generator-specific fingerprints rather than relying on a fixed feature
dictionary.

\subsection{CNN-Based Detectors}
\label{subsec:cnn_detectors}

Wang et al.~\cite{wang2020} established the dominant CNN-based detection paradigm with
their CNNSpot method. A ResNet~\cite{he2016} is fine-tuned on ProGAN-generated images
alongside real images from the LSUN dataset. At training time, Gaussian blur and JPEG
compression augmentation are applied to reduce reliance on high-frequency artefacts
specific to the training generator. The model is then evaluated on images from 11
different GAN architectures spanning a range of generator families.

The central finding of Wang et al.\ is that training on one generator produces near-perfect
in-distribution performance but substantially degraded performance on other generators.
Even with augmentation, cross-generator F1 drops to near-random for architectures whose
spectral footprint differs substantially from the training distribution. This finding,
that CNN detectors learn generator-specific shortcuts rather than universal artefact
signatures, is the central challenge addressed in this thesis. The work of Wang et al.\
was conducted entirely within the GAN paradigm on a single real-image source; this thesis
extends it to a cross-paradigm setting with three independent real-image sources and a
controlled four-stage training progression.

The ResNet architecture~\cite{he2016} is the backbone used in both CNNSpot and in all
14 experiments in this thesis. ResNet introduced residual connections, also called skip
connections, which allow a layer to learn the residual mapping $\mathcal{F}(x) = H(x) - x$
rather than the full transformation $H(x)$. This enables the training of much deeper
networks by providing a gradient path that bypasses each block and prevents vanishing
gradients. ResNet-50 (~25.6M parameters) and ResNet-18 (~11.7M parameters) are used in
this thesis, both initialised from ImageNet-pretrained weights and fine-tuned with a
modified two-class output layer.

\subsection{Universal and Architecture-Agnostic Detectors}
\label{subsec:universal_detectors}

Ojha et al.~\cite{ojha2023} proposed the most direct architectural response to the
cross-generator generalisation failure identified by Wang et al. Rather than fine-tuning
a CNN on synthetic images, they use a frozen CLIP ViT-L/14 backbone and classify images
using nearest-neighbour distance in CLIP feature space to a set of reference images.
Because the CLIP backbone is trained on a large and diverse corpus through a contrastive
text-image objective rather than on an in-distribution detection task, its feature
representations are less aligned with the statistical properties of any particular
generator. Tested on 21 generators including both GANs and diffusion models, the method
achieves substantially better cross-generator generalisation than CNN-based fine-tuning
on the same benchmarks.

The key insight of Ojha et al.\ is that the features learned by a task-specific
fine-tuned CNN are biased toward the training distribution in ways that do not transfer.
A backbone with broader pretraining produces a feature space where the boundary between
real and AI images is more stable across generators. This thesis studies the same failure
mode from a data composition angle rather than an architectural one: instead of changing
the feature extractor, the training corpus is diversified across multiple generators and
real-image sources. The two approaches are complementary, and the shortcut-learning failure
documented in the cross-dataset matrix provides a precise characterisation of what the
CLIP-based approach is designed to avoid.

\subsection{Spectral and Frequency Domain Approaches}
\label{subsec:spectral_detection}

Several detection approaches exploit the periodic spectral artefacts that GAN upsampling
grids produce. When a generator applies transposed convolution or nearest-neighbour
upsampling at a fixed stride, the operation introduces spectral peaks at frequencies
determined by that stride. These peaks appear in the FFT magnitude spectrum of generated
images as regular grid-like patterns absent from camera-captured photographs, whose
spectral content follows a natural $1/f$ decay. Detectors that operate in the frequency
domain or use spectral features as input can exploit these patterns directly.

A related preprocessing approach applies band-pass filtering, such as Difference of
Gaussians, to suppress low-frequency content and isolate the frequency bands where
generative artefacts concentrate. If the detector is supplied with a band-pass image
rather than the raw RGB input, the task of identifying spectral irregularities is
simplified. The DoG transform with parameters $\sigma_1 = 1.0$ and $\sigma_2 = 2.0$
produces a band-pass response centred roughly on the spatial frequency range
$[1/(2\sigma_2),\, 1/(2\sigma_1)]$ cycles per pixel, which covers the upsampling
artefact frequencies of typical GAN architectures with stride~2 or stride~4.

Diffusion models, which apply many small iterative refinements rather than a single
upsampling step, produce weaker and less consistent spectral artefacts than early GANs.
Each denoising step contributes a small correction that partially cancels the spectral
irregularities introduced by the previous step; the result is that modern diffusion
outputs lack the strong periodic peaks characteristic of transposed-convolution GANs.
Frequency-domain detection has therefore become less reliable as diffusion models have
displaced GANs as the dominant generation paradigm. Experiments~11 and~12 in Chapter~4
test the DoG preprocessing hypothesis directly on both diffusion-model and GAN datasets;
the results show consistent degradation relative to the RGB baseline and are analysed in
detail in Chapter~5.

% ─────────────────────────────────────────────────────────────────────────────

\section{Datasets for AI Image Detection}
\label{sec:detection_datasets}

A detection benchmark requires, at minimum, a balanced or known-imbalance split of real
and AI images, a held-out test set that withholds at least some generators from training,
and provenance metadata that identifies which generator produced each AI image. These
requirements are less commonly satisfied than in-distribution benchmarks, where a single
generator is used and the real-image source is fixed.

The ForenSynths benchmark introduced by Wang et al.~\cite{wang2020} comprises images
from 11 GAN architectures alongside real images from the LSUN large-scale scene dataset.
All 11 generators belong to the GAN paradigm. The real-image source is uniform across
all test subsets. ForenSynths was the primary benchmark for evaluating cross-generator
generalisation in the GAN era and established that in-distribution performance does not
predict cross-generator performance.

The UniversalFakeDetect benchmark introduced by Ojha et al.~\cite{ojha2023} extends the
generator coverage to 21 architectures, including both GAN and diffusion models. It
includes generators from both paradigms in the test set and is the benchmark on which
the CLIP-based universal detector was evaluated. Its inclusion of diffusion models makes
it the first benchmark to evaluate cross-paradigm generalisation systematically.

This thesis uses three datasets that together provide cross-paradigm, multi-generator,
and multi-real-source coverage. Dataset~A~\cite{rombach2022} is the ``AI vs.\
Human-Generated Images'' Kaggle dataset: approximately 80,000 balanced images with Stable
Diffusion~v1.5 as the AI source and Shutterstock stock photographs as the real source.
Dataset~B is the Defactify Image Dataset, available on HuggingFace: 96,000 images from
five diffusion-based generators (SD~2.1, SDXL, SD~3, DALL-E~3, MidJourney~v6) paired
with MS-COCO natural photographs at a 1:5 real-to-AI imbalance. Metadata identifies
which generator produced each AI image, enabling per-generator evaluation. Dataset~C~\cite{karras2019}
is the ``140k Real and Fake Faces'' Kaggle dataset: 140,000 balanced images from
StyleGAN~\cite{karras2019} paired with FFHQ portrait photographs.

None of the three datasets was designed for cross-dataset evaluation. Each was
constructed as an in-distribution benchmark for a specific generator or set of generators
against a specific real-image source. This thesis repurposes them together for a different
question: how well does a model trained on one dataset's joint distribution of real and AI
images transfer to another dataset's distribution? The answer depends not only on the
generator but also on the real-image source, a dependency that the cross-dataset matrix
makes explicit.

The choice of real-image source matters beyond academic convenience. Detectors deployed
in practice encounter images from sources that differ from their training data in both
generator and photographic subject matter. A model trained on studio portraits (FFHQ)
may never have learned the low-level statistics of natural-scene photographs (MS-COCO)
or stock images (Shutterstock); the resulting classifier may perform well on faces but
produce high false-positive rates on landscapes, products, or social media snapshots
where the real-image appearance differs substantially from training. Benchmark coverage
that varies only the generator while holding the real-image source constant cannot expose
this failure mode, because the real-class decision boundary is never tested outside its
training distribution.

% ─────────────────────────────────────────────────────────────────────────────

\section{Cross-Dataset Generalisation}
\label{sec:cross_dataset_generalisation}

The cross-dataset generalisation problem in AI image detection can be stated formally as
follows. Let $\mathcal{D}_k = \{(x_i, y_i)\}$ denote a dataset where each image $x_i$
is labelled $y_i \in \{\text{real}, \text{AI}\}$, with real images drawn from source
distribution $p_k^{\text{real}}$ and AI images from generator $g_k$. A detector trained
on $\mathcal{D}_{\text{train}}$ achieves in-distribution performance when evaluated on a
held-out test set from the same distribution, and out-of-distribution (OOD) performance
when evaluated on $\mathcal{D}_{\text{test}}$ where $p_{\text{test}}^{\text{real}} \neq
p_{\text{train}}^{\text{real}}$ or $g_{\text{test}} \neq g_{\text{train}}$, or both.

Wang et al.~\cite{wang2020} showed empirically that in-distribution performance does not
predict OOD performance: CNNSpot achieves over 0.99 accuracy on ProGAN test images but
collapses to near-random on images from other GAN architectures. The augmentation
strategy they propose improves OOD performance modestly but does not close the gap. Their
evaluation covers multiple generators but within a single paradigm (GANs) and a single
real-image source (LSUN), which means the two sources of distribution shift, generator
change and real-image source change, cannot be separated.

Ojha et al.~\cite{ojha2023} address the cross-generator problem architecturally, replacing
the fine-tuned CNN with a frozen CLIP backbone. Their results on 21 generators confirm
that a general-purpose feature space generalises better across generators than a
task-specific one. This approach targets the generator-shift component of OOD failure.

This thesis targets the same problem through training corpus design rather than feature
extractor design. The 3$\times$3 cross-dataset matrix covers both sources of distribution
shift simultaneously: generator change and real-image source change, across two generation
paradigms and three independent real-image sources. Prior benchmarks use a single
real-image source, so any OOD degradation is attributed entirely to the generator even
when the real-image shift is equally responsible; the matrix design here disentangles the
two by varying each independently. The four-stage progression then isolates the
contribution of each modification, establishing that real-image source diversity is the
primary driver of OOD improvement. Table~\ref{tab:method_comparison} summarises all three
approaches.

\begin{table}[htb]
  \centering
  \caption{Comparison of three detection approaches along key evaluation dimensions.}
  \label{tab:method_comparison}
  \resizebox{\linewidth}{!}{%
  \begin{tabular}{llll}
    \hline
    \textbf{Dimension} & \textbf{Wang~\cite{wang2020}} & \textbf{Ojha~\cite{ojha2023}} & \textbf{This thesis} \\
    \hline
    Feature extractor    & Fine-tuned ResNet    & Frozen CLIP ViT-L/14  & Fine-tuned ResNet \\
    Training paradigm    & Supervised fine-tune & Nearest-neighbour     & Supervised fine-tune \\
    Training generators  & 1 (ProGAN)           & 1 reference set       & 6 (diffusion + GAN) \\
    Test generators      & 11 (GAN only)        & 21 (GAN + diffusion)  & 6 (diffusion + GAN) \\
    Real-image sources   & 1 (LSUN)             & 1                     & 3 (Shutterstock, COCO, FFHQ) \\
    Cross-paradigm eval  & No                   & Yes                   & Yes \\
    Real-source ablation & No                   & No                    & Yes \\
    \hline
  \end{tabular}}
\end{table}
